% ===== PRÉAMBULE CANONIQUE — à coller VERBATIM en tête de CHAQUE .tex =====
\documentclass[11pt,a4paper]{article}
\usepackage[utf8]{inputenc}\usepackage[T1]{fontenc}\usepackage{lmodern}
\usepackage[french]{babel}
\usepackage{amsmath,amssymb,mathtools}\usepackage{icomma}
\usepackage{siunitx}\sisetup{locale=FR}  % \qty{}{}, \si{}, \ang{}
\usepackage{xcolor}\usepackage{colortbl}
\usepackage{tikz}\usepackage{pgfplots}\pgfplotsset{compat=1.18}
\usepackage[siunitx,europeanresistors]{circuitikz} % circuits (élec.)
\usepackage[most]{tcolorbox}\usepackage{enumitem}\usepackage{array,booktabs}
\usepackage[a4paper,margin=2.2cm]{geometry}
\usetikzlibrary{arrows.meta,calc,angles,quotes,positioning,patterns,%
  backgrounds,shapes.geometric,decorations.pathmorphing,decorations.markings,intersections}
\definecolor{primary}{RGB}{222,49,99}\definecolor{primarydark}{RGB}{150,22,55}
\definecolor{defcol}{RGB}{0,114,178}\definecolor{methcol}{RGB}{52,92,148}
\definecolor{excol}{RGB}{0,135,90}\definecolor{attncol}{RGB}{200,40,55}
\tcbset{boxbase/.style={enhanced,breakable,boxrule=0.9pt,arc=2.5pt,
  left=3mm,right=3mm,top=2.3mm,bottom=2.3mm,fonttitle=\bfseries,coltitle=white}}
% --- boîtes TUTORIEL (noms EXACTS, ne pas renommer) ---
\newtcolorbox{cle}[1][]{boxbase,colback=defcol!6,colframe=defcol,
  title={\textsc{À retenir}\ifx\relax#1\relax\relax\else\textmd{ : \itshape #1}\fi}}
\newtcolorbox{methode}[1][]{boxbase,colback=methcol!7,colframe=methcol,sharp corners,
  title={\textsc{Méthode}\ifx\relax#1\relax\relax\else\textmd{ --- \itshape #1}\fi}}
\newtcolorbox{exemplebox}[1][]{boxbase,colback=excol!5,colframe=excol,
  title={\textsc{Exemple}\ifx\relax#1\relax\relax\else\textmd{ : \itshape #1}\fi}}
\newtcolorbox{piege}[1][]{boxbase,colback=attncol!6,colframe=attncol,
  title={\textsc{Piège}\ifx\relax#1\relax\relax\else\textmd{ : \itshape #1}\fi}}
\newtcolorbox{remarque}[1][]{boxbase,colback=black!4,colframe=black!45,
  title={\textsc{Remarque}\ifx\relax#1\relax\relax\else\textmd{ : \itshape #1}\fi}}
% --- boîtes EXERCICE / CORRIGÉ (noms EXACTS, auto-comptées) ---
\newtcolorbox[auto counter]{exo}[1][]{boxbase,colback=primary!4,colframe=primary,
  title={\textsc{Exercice}~\thetcbcounter\ifx\relax#1\relax\relax\else\textmd{ --- \itshape #1}\fi}}
\newtcolorbox[auto counter]{corrige}[1][]{boxbase,colback=excol!4,colframe=excol,
  title={\textsc{Corrigé}~\thetcbcounter\ifx\relax#1\relax\relax\else\textmd{ --- \itshape #1}\fi}}
\newcommand{\vv}[1]{\vec{#1}}\newcommand{\ds}{\displaystyle}
\newcommand{\R}{\mathbb{R}}\newcommand{\N}{\mathbb{N}}\newcommand{\Z}{\mathbb{Z}}
% =========================================================================
\begin{document}
\usetikzlibrary{babel}
\begin{corrige}[intensité à partir de la charge]
On applique directement la définition de l'intensité $I=\dfrac{Q}{t}$ :
\[ I=\frac{Q}{t}=\frac{12}{4}=\qty{3}{\ampere}. \]
L'unité est bien l'ampère, car $\si{\ampere}=\si{\coulomb\per\second}$.
\end{corrige}

\begin{corrige}[de la charge au nombre d'électrons]
\begin{enumerate}[label=\alph*)]
  \item $Q=I\,t=2\times5=\qty{10}{\coulomb}$.
  \item $n=\dfrac{Q}{e}=\dfrac{10}{1{,}6\times10^{-19}}=6{,}25\times10^{19}$ électrons.
\end{enumerate}
\end{corrige}

\begin{corrige}[loi des nœuds : une intensité manquante]
Le courant $I$ arrive au nœud, les courants $I_1$ et $I_2$ en partent. La loi des nœuds donne
$I=I_1+I_2$, d'où :
\[ I_2=I-I_1=5-2=\qty{3}{\ampere}. \]
\end{corrige}

\begin{corrige}[en série, l'intensité est partout la même]
Le second ampèremètre lira aussi $\boxed{\qty{0.5}{\ampere}}$. En effet, $R_1$ et $R_2$ sont
\textbf{en série} : il n'y a aucun nœud entre eux, donc la même intensité les traverse d'un bout à
l'autre de la branche. La position de l'ampèremètre dans une branche série est sans importance.
\end{corrige}

\begin{corrige}[isoler la durée]
On part de $I=\dfrac{Q}{t}$, que l'on inverse en $t=\dfrac{Q}{I}$ :
\[ t=\frac{Q}{I}=\frac{30}{0{,}5}=\qty{60}{\second}\ (=\qty{1}{\minute}). \]
\end{corrige}

\end{document}
